\section{Timeline}
\label{sec:timeline}

The project is planned as a three-year effort, with thrusts serialized
in their primary focus but running in parallel at a lower intensity so
that an end-to-end prototype exists from the first year.

\paragraph{Year 1: Knowledge Base and Observation.}
Primary focus on Thrust~1. We will design and implement the hybrid
KB, develop the first generation of passive observation modules
(network, log, configuration, GitOps), and deploy them on the UCSB
internal network and on the healthcare-records reference testbed. By
the end of Year~1, a rudimentary Replication Agent and Red Team Agent
will operate against a hand-built twin of the healthcare-records
testbed, establishing the end-to-end loop.

\paragraph{Year 2: Replication.}
Primary focus on Thrust~2. We will formalize and implement the
controlled-fidelity replication algorithm, integrate direct
replication from upstream artifacts with Conware-inspired behavioral
mock synthesis for opaque components, and extend the Replication
Agent to consume the KB automatically. A second reference testbed
(the financial-services transaction-processing stack) and at least
one partner environment will be brought into the evaluation.
Year~2 deliverables include a quantitative study of twin fidelity on
the healthcare-records testbed.

\paragraph{Year 3: Analysis and integration.}
Primary focus on Thrust~3. We will develop the PDDL-based planning
domain for multi-host, multi-service open-source deployments,
implement the LLM/classical hybrid planner, integrate component-level
analyses as sources of atomic findings, and run end-to-end red-team
campaigns on all three reference testbeds. Year~3 will also include a
longitudinal evaluation of the feedback loop: how much does twin
fidelity and red-team precision improve over successive iterations of
analyst validation?

\paragraph{Cross-cutting activities.}
Open-source releases, publications, student involvement, and the
annual infrastructure-security-focused CTF competition
(\S\ref{sec:broader}) run across all three years. Each thrust
produces artifacts that feed the others, and the end-to-end prototype
is maintained continuously rather than assembled at the end.
